\chapter{Locations}\label{ch:in:locations}

A \$300\,000 package, after income tax and a year's rent of a two-bedroom flat by each city's official statistic, leaves
\$139\,220 in London and \$214\,938 in Hong Kong. The ordering is not fixed: New York leaves more than London at that
package and less above about \$851\,900, because the two tax schedules cross and the rents are almost equal. This chapter
describes the industry's cities, computes what a package is worth in five of them after tax and housing, sets out the
visas a move needs, and discusses what moving costs.

\section{The trading cities and what is there}

The industry's employers concentrate where its markets and clients are (chapter~1): New York and Chicago for US equities,
options and futures; London for currencies, rates and European markets; Amsterdam for European options and ETF market
making (chapters~2 and~3); Zurich for banks; Hong Kong and Singapore for Asia; and Dubai, whose absence of income tax
chapter~15 measured. Colocation (Book~14, chapter~9) ties the trading systems to the
exchanges' data centres, but not the people: a firm's researchers and engineers can sit far from its servers. The
chapter compares five cities with official rent statistics: London, New York, Chicago, Zurich and Hong Kong.

\begin{definition}[Price level index, disposable pay after housing]\label{def:in:locations:pli}
A \emph{price level index}\index{price level index} is the ratio of a country's purchasing power parity to its market
exchange rate, expressed against a reference area: above 100, the same basket costs more there than in the reference.
\emph{Disposable pay after housing}\index{disposable pay after housing} is a package's pay after income tax and employee
social contributions, minus a year's rent of a standard home in the city.
\end{definition}

\section{Pay, tax and housing}

\begin{dated}{2026-08}{Rents in the chapter's five cities}\label{dat:in:locations:rents}
London, August 2026: average private rent of a two-bedroom home in the London region, \pounds 2\,232 a month (ONS).
New York and Chicago, fiscal 2026: fair market rents of a two-bedroom unit, \$2\,910 and \$1\,781 a month; fair market
rents ``are estimates of 40th percentile gross rents for standard quality units'' (HUD). Zurich, March 2024: median net
rent of a three-room flat in the city, 1\,578 francs a month (City of Zurich). Hong Kong, August 2026 (provisional):
average rent of class~B flats (40 to 69.9 square metres) on Hong Kong Island, HK\$446 a square metre a month (Rating and
Valuation Department). Price level indices of actual individual consumption, 2024 (EU = 100): Switzerland 184.3, United
States 149.2, United Kingdom 129.1, Netherlands 121.0 (Eurostat).
\end{dated}

The five statistics measure different things: an average, a 40th percentile, a median; rents with utilities (HUD's
gross rents) and without (Zurich's net rent); a region (London), a metropolitan area (New York, Chicago) and a city
(Zurich). For Hong Kong the chapter multiplies the rent per square metre by an illustrative 55 square metres. The
comparison is therefore of orders of magnitude and of how the ranking moves, not of exact differences.

\begin{method}[Disposable pay after housing]\label{met:in:locations:disposable}
For a package $G$ in US dollars and a city $c$ with tax function $T_c$ (chapter~15's \texttt{firm.aftertax}, at the
ECB's 2025 average exchange rates) and monthly rent $r_c$ in local currency with $k_c$ units per dollar,
\[
D_c(G)=G-T_c(G)-12\,r_c/k_c .
\]
Two cities are compared by the package $G_b$ with $D_b(G_b)=D_a(G_a)$, and by the packages at which $D_a=D_b$.
\end{method}

\begin{omfigure}
\begin{tikzpicture}
\begin{axis}[xbar stacked,width=10.5cm,height=5cm,xmin=0,xmax=300,ymin=0.4,ymax=5.6,bar width=11pt,ytick=data,
  yticklabels from table={figdata/industry/27-locations/split.csv}{city},yticklabel style={font=\footnotesize},
  xticklabel style={font=\footnotesize},xlabel={\small \$ thousand of a \$300\,000 package},table/col sep=comma,
  legend style={legend cell align=left,font=\scriptsize,at={(0.5,-0.32)},anchor=north,legend columns=3,/tikz/every even column/.append style={column sep=8pt}}]
\addplot[fill=omDef,draw=none] table[x=disposable,y=pos]{figdata/industry/27-locations/split.csv};
\addplot[fill=omThm!60,draw=none] table[x=rent,y=pos]{figdata/industry/27-locations/split.csv};
\addplot[fill=black!25,draw=none] table[x=tax,y=pos]{figdata/industry/27-locations/split.csv};
\legend{disposable after housing,rent,tax and social contributions}
\end{axis}
\end{tikzpicture}

\omcaption{A \$300\,000 package split into \omterm{def:in:locations:pli}{disposable pay after housing}, a year's rent of a two-bedroom home by each city's
official statistic, and tax with employee social contributions (single employee, no other income, standard deductions).
Data: \texttt{data/industry/rents.csv}, chapter~15's tax data and two constants, through
\texttt{in\_locations.table}.}\label{fig:in:locations:split}
\end{omfigure}

At \$300\,000 (\cref{fig:in:locations:split}) the order is Hong Kong \$214\,938, Zurich \$184\,855, Chicago \$178\,955,
New York \$149\,062 and London \$139\,220. Tax separates the cities more than rent does: London's tax and contributions
take \$125\,455 and Hong Kong's \$47\,309, while the rents lie between \$21\,372 (Chicago) and \$37\,753 (Hong Kong). The
ordering moves with the package (\cref{fig:in:locations:vslondon}): New York is ahead of London from about \$134\,100 to
about \$851\,900 and behind outside that range, because New York's schedule taxes large incomes more heavily than
London's (chapter~15) while their rents are within \$400 a year of each other.

\begin{omfigure}
\begin{tikzpicture}
\begin{axis}[width=10.5cm,height=5.8cm,xmin=100,xmax=1000,ymin=-15,ymax=80,
  xlabel={\small package, \$ thousand},ylabel={\small disposable pay minus London's, \$ thousand},
  tick label style={font=\footnotesize},grid=major,grid style={gray!20},table/col sep=comma,
  legend style={font=\scriptsize,at={(0.97,0.22)},anchor=south east,legend cell align=left}]
\addplot[omThm,thick] table[x=gross,y=Hong_Kong]{figdata/industry/27-locations/vs_london.csv};
\addplot[black!70,thick,dashed] table[x=gross,y=Zurich]{figdata/industry/27-locations/vs_london.csv};
\addplot[omDef!60,thick] table[x=gross,y=Chicago]{figdata/industry/27-locations/vs_london.csv};
\addplot[omDef,thick] table[x=gross,y=New_York]{figdata/industry/27-locations/vs_london.csv};
\draw[black!50] (axis cs:100,0) -- (axis cs:1000,0);
\legend{Hong Kong,Zurich,Chicago,New York}
\end{axis}
\end{tikzpicture}

\omcaption{\omterm{def:in:locations:pli}{Disposable pay after housing} in four cities minus London's, for packages from \$100\,000 to \$1 million. The
New York line crosses zero twice; the Hong Kong line leaves the chart just above \$300\,000 and reaches \$296\,131 at \$1
million. Data: \texttt{in\_locations.curves}.}\label{fig:in:locations:vslondon}
\end{omfigure}

Price levels add a third layer that the chapter does not model: the same disposable pay buys less in Switzerland,
whose price level for household consumption was 184.3 against the European Union's 100 in 2024, than in the United
Kingdom (129.1) or the Netherlands (121.0) (\cref{fig:in:locations:pli}). The chapter subtracts rent because official
statistics measure it city by city; the rest of the cost of living is better handled by such an index than by
city-specific guesses.

\begin{omfigure}
\begin{tikzpicture}
\begin{axis}[xbar,width=10.5cm,height=4cm,xmin=0,xmax=210,ymin=0.4,ymax=4.6,bar width=9pt,ytick=data,
  yticklabels from table={figdata/industry/27-locations/pli.csv}{country},yticklabel style={font=\footnotesize},
  xticklabel style={font=\footnotesize},xlabel={\small price level index, actual individual consumption, EU = 100, 2024},
  grid=major,grid style={gray!20},table/col sep=comma,nodes near coords,
  nodes near coords style={font=\scriptsize,anchor=west,xshift=1pt},point meta=explicit symbolic]
\addplot[fill=omDef,draw=none] table[x=pli,y=pos,meta=pli]{figdata/industry/27-locations/pli.csv};
\end{axis}
\end{tikzpicture}

\omcaption{Price level indices of household consumption in four countries, 2024, from Eurostat's purchasing power
parities. Data: Eurostat \texttt{prc\_ppp\_ind}, through \texttt{in\_locations.PLI\_2024}.}\label{fig:in:locations:pli}
\end{omfigure}

\section{Visas and work permits}

\begin{definition}[Specialty-occupation visa, skilled-worker visa]\label{def:in:locations:visas}
A \emph{specialty-occupation visa}\index{specialty-occupation visa} is a US temporary work status (H-1B) for jobs
requiring at least a bachelor's degree in a specific field, sponsored by an employer who files a labour condition
application (chapter~14). A \emph{skilled-worker visa}\index{skilled-worker visa} is the United Kingdom's sponsored work
route, for a job at a required skill level paid at least a salary threshold.
\end{definition}

\begin{dated}{2026-09}{Visas for the industry's cities}\label{dat:in:locations:visas}
United States (USCIS, September 2026): H-1B regular cap 65\,000 a year plus 20\,000 for holders of US master's degrees;
from fiscal year 2027 a ``weighted selection process which favors the allocation of H-1B visas to higher skilled and
higher paid'' applicants; a proclamation of 19 September 2025 requires an additional \$100\,000 payment for certain
petitions filed from 21 September 2025, whose implementing guidance a district court vacated on 8 June 2026, an order
administratively stayed while the government appeals. United Kingdom: Skilled Worker visa at ``at least \pounds 41\,700
per year, or the 'going rate' for your job, whichever is higher''. Singapore: Employment Pass qualifying salary for
financial services S\$6\,200 a month, rising with age to S\$11\,800, and S\$6\,600 for new applications from 1 January
2027, plus the COMPASS points test. Netherlands, 2026: highly skilled migrant or EU Blue Card, EUR 5\,942 a month gross
(30 or older), EUR 4\,357 (under 30).
\end{dated}

For a quantitative hire the thresholds are rarely binding (chapter~14's offers sit far above them); the constraints are
elsewhere. The US lottery decides whether a first H-1B is possible at all in a given year, and the 2027 weighting and the
disputed \$100\,000 payment change the odds and the cost for the employer. Other countries' routes are closer to an
administrative check, which is one reason firms open offices where hiring from abroad is predictable.

\section{Moving: relocation, remote work and the cost of leaving}

A move has three costs beyond the new city's tax and rent. \textbf{Relocation}: flights, temporary housing, deposits and
the months of double rent, which an employer may cover for a hire. \textbf{The cost of leaving}: unvested deferred pay and
non-compete periods at the old employer (chapters~13 and~28), which the new employer may buy out. \textbf{\omterm{def:in:pay-regulation-and-tax-by-location:expat}{Tax
residence}}: moving mid-year can split a year between two tax systems (chapter~15). Remote work changes the first but not the last: an employee who works from another country
may create tax and regulatory obligations for the employer, one reason a firm whose staff handle regulated activities
may require them to work from an office where it is licensed.

\section{Tutorial: the same offer in five cities}\label{tut:in:locations:tut}

\begin{tutorial}
\textbf{Goal.} Compute \omterm{def:in:locations:pli}{disposable pay after housing} for a package in five cities, and the packages at which two cities are
equal.
\textbf{End state:} \cref{fig:in:locations:split,fig:in:locations:vslondon,fig:in:locations:pli}.
\begin{tutsteps}
\item \textbf{Cities.} \texttt{firm.locations.City} records join a \texttt{firm.aftertax} location, a rent with its
measure and period, and a ledger id; \texttt{data/industry/rents.csv} holds the ONS and HUD rows.
\item \textbf{Disposable pay.} \texttt{disposable(params, city, gross)} subtracts a year's rent from the net pay
(\cref{lst:in:locations:disposable}).
\omcode{code/firm/locations/firm_locations.py}{35}{42}{A year's rent in dollars and disposable pay after housing.}\label{lst:in:locations:disposable}
\item \textbf{Comparisons.} \texttt{equivalent} finds the package in one city that matches another city's disposable
pay; \texttt{crossing} the package at which two cities are equal, by bisection.
\end{tutsteps}
\end{tutorial}

For \$300\,000 in London the equivalent New York package is \$281\,047, and for \$300\,000 in New York the equivalent
London package \$318\,569; the two cities cross at about \$134\,100 and \$851\,900.

\textbf{What to change next.} Replace the rent statistics with one measure for all cities (a median two-bedroom asking rent
from one provider, with its licence); add the bonus-deferral profile (chapter~13), which changes the year of taxation;
add schooling and commuting for a family.

\section{Build: firm.locations}\label{bld:in:locations:locations}

\begin{build}
\bfield{Purpose} Compare what a package leaves after tax and housing in different cities, and find the packages that make
two cities equal.
\bfield{Interface} \texttt{firm.locations}: \texttt{City(name, location, rent\_month, currency, measure, period, source,
visas)}; \texttt{rent\_usd\_year}; \texttt{disposable(params, city, gross)}; \texttt{equivalent(params, a, gross\_a, b)};
\texttt{crossing(params, a, b, lo, hi)}; on \texttt{firm.aftertax}.
\bfield{Rules} Every rent carries its measure, period and source; the exchange rates are the tax data's; a crossing is
reported only inside the searched range, and the caller chooses the range.
\bfield{Acceptance tests} \texttt{code/firm/locations/tests/}: with no tax, disposable pay is gross minus rent;
\texttt{equivalent} inverts \texttt{disposable}; two cities that never cross give no crossing.
\bfield{Stretch} Price level adjustment; families (two earners, children); a Monte Carlo over the bonus.
\end{build}

\begin{omsources}
\item ONS, Price Index of Private Rents; HUD, Fair Market Rents FY2026; City of Zurich, Mietpreiserhebung 2024; Hong Kong
Rating and Valuation Department; Eurostat, purchasing power parities.
\item USCIS, H-1B pages; GOV.UK, Skilled Worker visa; Singapore MOM, Employment Pass; IND, required amounts 2026.
\item Chapter~15's tax data and \texttt{firm.aftertax}.
\end{omsources}

\section{Exercises}

\begin{exercise}[$\star$]\label{exo:in:locations:1}
What is a year's rent in dollars in London, and in Hong Kong, by the chapter's statistics?
\end{exercise}

\begin{exercise}[$\star$]\label{exo:in:locations:2}
By how much does Hong Kong's disposable pay exceed London's at \$300\,000, and what share of the difference is tax?
\end{exercise}

\begin{exercise}[$\star$]\label{exo:in:locations:3}
What monthly salary must a 30-year-old highly skilled migrant earn in the Netherlands in 2026, and what is that a year?
\end{exercise}

\begin{exercise}[$\star\star$]\label{exo:in:locations:4}
Why do New York and London cross twice?
\end{exercise}

\begin{exercise}[$\star\star$]\label{exo:in:locations:5}
Zurich's rent is a 2024 median net of charges; London's is a 2026 average. In which direction does each difference bias
the comparison?
\end{exercise}

\begin{exercise}[$\star\star$]\label{exo:in:locations:6}
A candidate compares \$300\,000 in New York with \$280\,000 in Chicago. Which leaves more after tax and rent, and what
else should she weigh?
\end{exercise}

\begin{exercise}[$\star\star\star$]\label{exo:in:locations:7}
\emph{Coding.} Double the Hong Kong flat to 110 square metres. What happens to Hong Kong's disposable pay at \$300\,000
and to its lead over Zurich?
\end{exercise}

\begin{exercise}[$\star\star\star$]\label{exo:in:locations:8}
\emph{Find the flaw.} ``Zurich pays the most after tax and rent among the European cities, so it is the best place to
live on a quant's salary.''
\end{exercise}

\section{Problem: The Same Offer in Five Cities}

\begin{problem}[{Weekend problem --- the same offer in five cities}]\label{pb:in:locations:1}
A researcher has the same \$300\,000 offer from one firm's offices in London, New York, Chicago, Zurich and Hong Kong.

\medskip\noindent\textbf{Part I --- The cities.}
\begin{enumerate}
\item Which of the industry's cities does the chapter compare, and why these?
\item Define the \omterm{def:in:locations:pli}{price level index} and \omterm{def:in:locations:pli}{disposable pay after housing}.
\item State the five rent statistics and how they differ.
\item Why can colocation not decide where people work?
\item Which measures would make the comparison more consistent?
\end{enumerate}

\medskip\noindent\textbf{Part II --- The numbers.}
\begin{enumerate}[resume]
\item State the method.
\item Give the net pay, rent and disposable pay in each city at \$300\,000.
\item Which separates the cities more, tax or rent?
\item Give the equivalent packages between London and New York.
\item At which packages do London and New York cross?
\end{enumerate}

\medskip\noindent\textbf{Part III --- Visas and moving.}
\begin{enumerate}[resume]
\item Define the specialty-occupation and \omterm{def:in:locations:visas}{skilled-worker visas}.
\item State the H-1B cap, the weighting from fiscal 2027 and the status of the \$100\,000 payment.
\item State the UK, Singapore and Dutch thresholds.
\item What are the three costs of moving beyond tax and rent?
\item Why do trading firms tend to require presence in an office?
\end{enumerate}

\medskip\noindent\textbf{Part IV --- The verdict.}
\begin{enumerate}[resume]
\item State the \emph{named result}: disposable pay after tax and a year's two-bedroom rent for the chapter's package in
five cities, and the gross package that equalises London and New York.
\item How would price levels change the ranking?
\item What would change at \$1 million?
\item What does the visa situation add for a non-US citizen choosing New York?
\item In two sentences, advise the researcher.
\end{enumerate}
\end{problem}

\section{Interview questions}

\begin{interviewq}[$\star$ \iqroles{researcher}]\label{iq:in:locations:1}
How would you compare two offers in different countries?
\end{interviewq}

\begin{interviewq}[$\star$ \iqroles{developer}]\label{iq:in:locations:2}
Why might a firm keep its engineers in the city of its main exchange rather than anywhere?
\end{interviewq}

\begin{interviewq}[$\star\star$ \iqroles{researcher}]\label{iq:in:locations:3}
Two cities' \omterm{def:in:pay-regulation-and-tax-by-location:rates}{average tax rates} are equal at \$300\,000. Can their \omterm{def:in:locations:pli}{disposable pay after housing} differ by more than the rent
difference?
\end{interviewq}

\begin{interviewq}[$\star\star$ \iqroles{researcher, developer}]\label{iq:in:locations:4}
Design a function that returns the break-even package between two cities, and say what can go wrong numerically.
\end{interviewq}

\begin{interviewq}[$\star\star$ \iqroles{trader}]\label{iq:in:locations:5}
Why might a market maker keep offices in several time zones rather than in one place?
\end{interviewq}

\begin{interviewq}[$\star\star\star$ \iqroles{researcher}]\label{iq:in:locations:6}
The dollar falls 10\% against the pound. What happens to the London side of the chapter's comparison, and to a London
employee paid in pounds?
\end{interviewq}
